\documentclass[11pt]{article}
\usepackage[T1]{fontenc}
\usepackage{textcomp}
\usepackage[margin=0.75in]{geometry}
\usepackage{amsmath,amssymb,amsthm}
\usepackage{array,booktabs}
\usepackage{hyperref}
\setlength{\parindent}{0pt}
\newtheorem{theorem}{Theorem}
\theoremstyle{definition}
\newtheorem{definition}{Definition}
\title{Flow Proof Artifact: prop-09-similar-subtriangles-imply-right}
\author{Generated by \texttt{flow doc proof}}
\date{Flow Proof Book}
\begin{document}
\maketitle
If from a point on the base of a triangle a line is drawn meeting the sides and the subtriangles are similar to the whole, the angle at the base is right.\\[0.5em]
\textbf{Source.} euclid — Elements, Book VI, Proposition 9\\[0.5em]
\bigskip
\noindent{\large \textbf{Derived fact 1.} \textit{If from a point on the base of a triangle a line is drawn meeting the sides and the subtriangles are similar to the whole, the angle at the base is right} \hfill the Euclidean plane \cdot Euclid Book VI \cdot Proposition 9: similar subtriangles sharing an acute angle imply a right triangle}\par
\smallskip\noindent\textit{Coordinate: the Euclidean plane \cdot Euclid Book VI \cdot Proposition 9: similar subtriangles sharing an acute angle imply a right triangle}\par
\smallskip\noindent\textit{Source: euclid — Elements, Book VI, Proposition 9}\par
\smallskip\noindent\textit{Built on: proposition 8: in a right triangle the altitude to the hypotenuse gives similar subtriangles, for Euclid Book VI on the Euclidean plane}\par
\medskip
\noindent\textbf{Goal.} If from a point on the base of a triangle a line is drawn meeting the sides and the subtriangles are similar to the whole, the angle at the base is right\par
\noindent$\displaystyle \text{similar subtriangles implies right angle}$\par
\medskip
\noindent
\begin{tabular}{@{} >{\raggedright\arraybackslash}p{0.06\textwidth} >{\raggedright\arraybackslash}p{0.47\textwidth} >{\raggedleft\arraybackslash}p{0.41\textwidth} @{}}
\textbf{\#} & \textbf{Proof} & \textbf{Mathematics} \\
\midrule
\textbf{1.} & We prove that proposition 9: similar subtriangles sharing an acute angle imply a right triangle for Euclid Book VI the Euclidean plane. & \\[0.45em]
\textbf{2.} & Let ABC = triangle. & \\[0.45em]
\textbf{3.} & Let D = point\_on\_base. & \\[0.45em]
\textbf{4.} & Let subtriangle = ABD. & \\[0.45em]
\textbf{5.} & From step 2, step 3, and step 4, we can deduce that angle BAC equals one right angle. & $\displaystyle \angle BAC = 90^\circ$ \\[0.45em]
\textbf{6.} & We invoke the derived fact governing Euclid Book VI the Euclidean plane: proposition 8: in a right triangle the altitude to the hypotenuse gives similar subtriangles, for Euclid Book VI on the Euclidean plane. & \\[0.45em]
\textbf{7.} & From step 6, this implies similar subtriangles implies right angle. Hence proven. & $\displaystyle \text{similar subtriangles implies right angle}$ \\[0.45em]
\bottomrule
\end{tabular}
\label{thm:Geometry-Euclid Book VI-proposition_9_similar_subtriangles_sharing_an_acute_angle_imply_a_right_triangle}
\par\medskip\hrule
\end{document}
